\section{Elemen Semantik dan Struktur Halaman}

HTML5 memperkenalkan serangkaian elemen semantik baru yang merevolusi cara pengembang menandai struktur halaman web. Sebelum HTML5, pengembang sering menggunakan elemen \code{<div>} dengan atribut \code{id} atau \code{class} seperti \code{<div id="header">} atau \code{<div class="nav">} untuk membedakan bagian halaman. Pendekatan ini berfungsi untuk styling CSS namun tidak memberikan makna struktural bagi browser, mesin pencari, atau teknologi bantu. Elemen semantik HTML5 seperti \code{<header>}, \code{<nav>}, \code{<main>}, \code{<article>}, \code{<section>}, \code{<aside>}, dan \code{<footer>} menggantikan pola \code{<div>} generik dengan tag yang secara inheren menjelaskan makna dan peran konten di dalamnya \cite{webdevHtml}.

Elemen \code{<header>} merepresentasikan bagian pengantar konten atau bagian navigasi. Ini dapat digunakan sebagai header seluruh halaman atau sebagai header untuk elemen tertentu seperti \code{<article>} atau \code{<section>}. Elemen \code{<nav>} menandakan blok link navigasi mayor - biasanya menu utama website. Elemen \code{<main>} adalah tempat konten utama halaman yang unik dan spesifik untuk dokumen tersebut; setiap halaman hanya boleh memiliki satu elemen \code{<main>} dan tidak boleh menempatkannya di dalam elemen struktural lain seperti \code{<article>} atau \code{<header>}. Elemen \code{<article>} merujuk pada konten mandiri yang dapat berdiri sendiri dan didistribusikan secara independen, seperti posting blog, artikel berita, atau kartu produk \cite{mdnAccessibility}.

Elemen \code{<section>} merepresentasikan bagian thematis dari dokumen, biasanya dengan heading sendiri. Elemen ini cocok untuk mengelompokkan konten terkait dalam satu bagian yang logis. Elemen \code{<aside>} berisi konten yang berkaitan secara tangensial dengan konten utama - sidebar, call-out boxes, atau iklan yang melengkapi tetapi terpisah dari alur utama. Elemen \code{<footer>} berisi informasi penutup untuk konten induknya atau untuk seluruh halaman, seperti informasi hak cipta, link legal, atau informasi kontak. Penggunaan semantik yang benar memungkinkan pembaca layar untuk menavigasi halaman dengan menu landmark, melompat langsung ke navigasi, konten utama, atau footer.

Manfaat menggunakan elemen semantik meluas di berbagai aspek pengembangan web. Untuk \textbf{SEO}, mesin pencari menggunakan struktur semantik untuk lebih memahami hierarki dan konteks konten, berpotensi meningkatkan peringkat pencarian. Untuk \textbf{aksesibilitas}, teknologi bantu seperti pembaca layar dapat menggunakan landmark semantik untuk memberikan navigasi yang lebih efisien kepada pengguna dengan disabilitas. Untuk \textbf{maintainability}, kode dengan elemen semantik lebih mudah dibaca dan dipelihara oleh tim pengembang karena struktur halaman jelas tanpa harus memeriksa atribut class. Untuk \textbf{konsistensi}, struktur semantik yang standar memudahkan integrasi dengan framework dan sistem manajemen konten.

\begin{table}[htbp]
\centering
\begin{tabular}{|P{3cm}|P{7cm}|}
\hline
\textbf{Elemen} & \textbf{Fungsi dan Penggunaan} \\
\hline
\code{<header>} & Bagian pengantar: logo, judul, navigasi awal \\
\hline
\code{<nav>} & Blok link navigasi mayor \\
\hline
\code{<main>} & Konten utama halaman (hanya satu per halaman) \\
\hline
\code{<article>} & Konten mandiri yang dapat didistribusikan (blog post, berita) \\
\hline
\code{<section>} & Bagian thematis dengan heading \\
\hline
\code{<aside>} & Konten sampingan terkait (sidebar, call-out) \\
\hline
\code{<footer>} & Informasi penutup: copyright, link legal \\
\hline
\end{tabular}
\caption{Elemen Semantik Struktur Halaman HTML5}
\label{tab:semantic-elements}
\end{table}

\begin{figure}[htbp]
\centering
\begin{tikzpicture}[
  node distance=0.5cm,
  every node/.style={font=\small},
  header/.style={rectangle, draw, fill=blue!20, minimum width=10cm, minimum height=0.8cm, align=center},
  nav/.style={rectangle, draw, fill=green!20, minimum width=10cm, minimum height=0.6cm, align=center},
  main/.style={rectangle, draw, fill=orange!20, minimum width=6.5cm, minimum height=3cm, align=center},
  aside/.style={rectangle, draw, fill=yellow!20, minimum width=3cm, minimum height=3cm, align=center},
  footer/.style={rectangle, draw, fill=gray!30, minimum width=10cm, minimum height=0.8cm, align=center},
  article/.style={rectangle, draw, fill=orange!30, minimum width=5.5cm, minimum height=1cm, align=center},
  section/.style={rectangle, draw, fill=orange!25, minimum width=5.5cm, minimum height=1cm, align=center}
]

\node[header] (header) {\code{<header>} - Logo, Judul};
\node[nav, below=of header] (nav) {\code{<nav>} - Menu Navigasi};
\node[main, below=of nav, xshift=-1.7cm] (main) {};
\node[aside, below=of nav, xshift=3.3cm] (aside) {\code{<aside>} - Sidebar};
\node[article, below=0.3cm of nav.north, anchor=north, xshift=-1.7cm] (article1) {\code{<article>} - Konten 1};
\node[section, below=0.2cm of article1] (section1) {\code{<section>} - Bagian A};
\node[footer, below=of main.south, yshift=-0.5cm] (footer) {\code{<footer>} - Copyright, Link};

\node at (main.center) [yshift=0.3cm] {\code{<main>}};
\node at (main.center) [yshift=-0.3cm, font=\footnotesize] {Konten Utama};

\end{tikzpicture}
\caption{Struktur Halaman Tipikal dengan Elemen Semantik HTML5}
\label{fig:semantic-structure}
\end{figure}

\begin{webcode}[language=HTML, caption={Contoh Struktur Halaman dengan Elemen Semantik}]
<!DOCTYPE html>
<html lang="id">
<head>
  <title>Contoh Struktur Semantik</title>
</head>
<body>
  <header>
    <h1>Nama Website</h1>
    <nav>
      <ul>
        <li><a href="/">Beranda</a></li>
        <li><a href="/tentang">Tentang</a></li>
        <li><a href="/kontak">Kontak</a></li>
      </ul>
    </nav>
  </header>
  
  <main>
    <article>
      <header>
        <h2>Judul Artikel</h2>
        <p>Dipublikasikan: 3 Maret 2026</p>
      </header>
      <section>
        <h3>Pendahuluan</h3>
        <p>Isi pendahuluan artikel...</p>
      </section>
      <section>
        <h3>Pembahasan</h3>
        <p>Isi pembahasan...</p>
      </section>
    </article>
  </main>
  
  <aside>
    <h2>Artikel Terkait</h2>
    <ul>
      <li><a href="#">Link 1</a></li>
      <li><a href="#">Link 2</a></li>
    </ul>
  </aside>
  
  <footer>
    <p>&copy; 2026 Nama Website. Hak Cipta Dilindungi.</p>
  </footer>
</body>
</html>
\end{webcode}

